\chapter{Introducción}
\label{chap:introduccion}

\section{Contexto}
\label{sec:contexto}

La gestión de proyectos es un proceso complejo que requiere la coordinación de múltiples tareas, recursos y personas. Tradicionalmente, las herramientas de gestión han utilizado representaciones tabulares y listas para organizar la información, lo que puede resultar en una sobrecarga cognitiva para los usuarios.

En los últimos años se ha innovado en el campo de la visualización de datos en diversos sectores, como los videojuegos, la simulación, la monitorización de sistemas y la visualización científica. Estas áreas han desarrollado técnicas avanzadas para representar información compleja de manera intuitiva, esto sin embargo no se ha trasladado de manera efectiva al ámbito de la gestión de proyectos.

Por otro lado, la inteligencia artificial es el foco de investigación, desarrollo y aplicación en los últimos años, y ha demostrado ser capaz de automatizar tareas repetitivas, mejorar la toma de decisiones y optimizar procesos en diversos campos de forma efectiva.

\section{Motivación}
\label{sec:motivacion_objetivos}

Aunque las herramientas actuales de gestión han integrado con éxito visualizaciones clásicas, estas en su mayoría perpetúan paradigmas analógicos. Esta herencia visual desaprovecha las capacidades de las interfaces modernas para sintetizar la complejidad, lo que genera dos barreras principales en la industria actual:

\begin{itemize}
    \item \textbf{Alta carga cognitiva y fatiga visual}: Para comprender el estado real de un proyecto, los gestores y desarrolladores deben procesar manualmente una gran cantidad de texto y datos abstractos dispersos. La falta de representaciones visuales intuitivas dificulta la detección rápida de bloqueos, dependencias o riesgos, obligando al usuario a realizar un esfuerzo mental considerable para construir un modelo mental del proyecto.
    \item \textbf{Burocracia frente a trabajo de valor}: Gran parte del tiempo de gestión se invierte en tareas mecánicas y repetitivas: solicitar actualizaciones de estado, requerir la cumplimentación de campos y sintetizar información fragmentada. Esta microgestión interrumpe los estados de concentración de los perfiles técnicos y desvía a los responsables de la toma de decisiones estratégicas.
\end{itemize}

\section{Objetivos}

\subsection{Objetivo General}
Investigar, diseñar y desarrollar un nuevo paradigma de plataforma de gestión de proyectos que reemplace las vistas tabulares convencionales por metáforas visuales inmersivas, capaces de representar la complejidad estructural del trabajo. Asimismo, se pretende concebir e integrar un sistema inteligente capaz de actuar de forma autónoma para asistir en la coordinación asíncrona y el seguimiento de los equipos.

\subsection{Objetivos Específicos}
Para alcanzar el objetivo general, se plantean los siguientes objetivos específicos:

\begin{itemize}
    \item \textbf{A. Investigación y análisis}: Analizar las limitaciones de usabilidad y \ac{UX} en las herramientas de gestión convencionales y estudiar su impacto en la carga cognitiva. Investigar patrones de diseño de interfaces en sectores ajenos a la gestión empresarial que permitan representar jerarquías y dependencias de forma intuitiva.
    
    \item \textbf{B. Diseño conceptual y visual}: Definir una propuesta de interfaz gráfica y experiencia de usuario que permita visualizar la topología de un proyecto como un sistema interconectado. El objetivo es priorizar la claridad visual frente a la lectura extensiva de texto, fomentando una interacción más natural y adaptada a entornos profesionales.
    
    \item \textbf{C. Sistema de agentes}: Diseñar un sistema basado en \ac{IA} capaz de actuar como intermediario ágil entre el usuario y la plataforma. El propósito es establecer flujos de trabajo que permitan delegar tareas repetitivas de coordinación, comunicación asíncrona y seguimiento de estados.
    
    \item \textbf{D. Arquitectura de software}: Diseñar e implementar una base técnica robusta que garantice la escalabilidad, el alto rendimiento y la mantenibilidad del código. Se persigue estructurar el sistema aplicando principios de diseño de software avanzados que aseguren un alto grado de desacoplamiento entre los dominios de la aplicación, facilitando su evolución y mantenimiento futuros.
    
    \item \textbf{E. Prototipado y validación técnica}: Desarrollar un \ac{MVP} que integre los nuevos paradigmas de visualización y el sistema multi-agente. Validar la viabilidad de la arquitectura propuesta y evaluar si la solución aporta un valor funcional tangible frente a los enfoques de gestión de proyectos tradicionales.
\end{itemize}

\section{Estado del Arte}

La gestión de proyectos de software es un dominio maduro con multitud de herramientas consolidadas. No obstante, la constante evolución de las metodologías de trabajo y el auge de la \ac{IA} están forzando un cambio de paradigma en cómo se visualiza la información y cómo se interactúa con el sistema.

\subsection{Paradigmas de visualización}
Las herramientas líderes en el sector, como Jira, Asana o Trello, han evolucionado significativamente ofreciendo múltiples perspectivas de los datos. Tradicionalmente estructuradas en representaciones tabulares y listas planas, la mayoría ha incorporado con éxito tableros Kanban para visualizar el flujo de trabajo y diagramas de Gantt para la planificación temporal.

Si bien los tableros Kanban son excelentes para gestionar el estado de tareas aisladas, presentan dificultades a la hora de representar relaciones jerárquicas complejas y dependencias cruzadas en proyectos densos. Por su parte, los diagramas de Gantt ofrecen una visión cronológica precisa, pero tienden a volverse visualmente sobrecargados y rígidos, resultando a menudo poco ágiles para el trabajo diario de los perfiles técnicos. Esto genera una fragmentación donde el usuario debe cambiar constantemente de vista para construir un modelo mental completo del proyecto, aumentando la carga cognitiva. El ecosistema actual evidencia un área de mejora en la adopción de representaciones topológicas interactivas (como grafos o vistas de árbol bidimensionales) que permitan observar simultáneamente la jerarquía, el estado y las dependencias estructurales de forma orgánica y sin depender exclusivamente del eje temporal.

\subsection{Evolución de la Inteligencia Artificial}
La integración de \ac{LLM} en la gestión de proyectos está experimentando una rápida transición. Las primeras implementaciones, visibles en herramientas como Atlassian Intelligence o Notion AI, se han centrado en un enfoque reactivo: asistentes conversacionales diseñados para resumir hilos de comentarios, redactar descripciones o traducir requisitos, actuando siempre bajo demanda síncrona del usuario.

Sin embargo, la tendencia del mercado avanza hacia la autonomía funcional. Plataformas como ClickUp están comenzando a explorar flujos de trabajo basados en agentes, donde la \ac{IA} puede encadenar acciones simples y automatizaciones. A pesar de estos avances, la mayoría de estas soluciones todavía operan como capas externas o complementos sobre sistemas heredados. El reto técnico y funcional de la industria reside actualmente en la integración nativa de estos agentes como actores de primer nivel: entidades proactivas capaces no solo de ejecutar herramientas internas, sino de gestionar su propio ciclo de vida, participando en los canales de comunicación asíncrona y suspendiendo sus ejecuciones a la espera de validación humana cuando la criticidad de la acción lo requiere.